\documentclass[conference]{IEEEtran}
\IEEEoverridecommandlockouts

% ---------------------------------------------------------
% PAQUETES
% ---------------------------------------------------------
\usepackage[utf8]{inputenc}
\usepackage[spanish,es-tabla]{babel}
\usepackage{amsmath,amssymb}
\usepackage{graphicx}
\usepackage{booktabs}
\usepackage{multirow}
\usepackage{array}
\usepackage{longtable}
\usepackage{hyperref}
\usepackage{float}
\usepackage{xcolor}
\usepackage{enumitem}

\hypersetup{
    colorlinks=true,
    linkcolor=black,
    citecolor=black,
    urlcolor=blue
}

% ---------------------------------------------------------
% INFORMACIÓN DEL DOCUMENTO
% ---------------------------------------------------------

\title{Comportamiento Adaptativo y Toma de Decisiones en Agentes Ecológicos mediante Aprendizaje por Refuerzo: Una Revisión de Literatura}

\author{
\IEEEauthorblockN{Axel López Vega}
\and
\IEEEauthorblockN{Kevin Núñez Camacho}
\and
\IEEEauthorblockN{Felipe Murillo}
}

% ---------------------------------------------------------
% DOCUMENTO
% ---------------------------------------------------------

\begin{document}

\maketitle

% ---------------------------------------------------------
% RESUMEN
% ---------------------------------------------------------

\begin{abstract}

El estudio computacional del comportamiento animal requiere representar
individuos capaces de responder ante cambios en el ambiente, disponibilidad
de recursos, amenazas y relaciones con otros individuos. Los Modelos
Basados en Agentes (ABM) permiten representar estos sistemas mediante
entidades autónomas, mientras que el Aprendizaje por Refuerzo (RL) permite
estudiar cómo dichas entidades pueden aprender estrategias de decisión a
partir de su interacción con el entorno.

Este trabajo presenta una revisión de literatura de 21 trabajos
relacionados con aprendizaje por refuerzo, modelado basado en agentes,
teoría de forrajeo, sistemas multiagente y comportamiento emergente. La
revisión se orienta mediante dos preguntas de investigación relacionadas
con la emergencia de estrategias adaptativas ante escenarios de estrés o
escasez de recursos y con la selección de variables de estado y funciones
de recompensa para representar decisiones ecológicas.

Los trabajos revisados permiten identificar como elementos recurrentes la
incorporación de variables fisiológicas, información espacial, percepción
parcial, memoria y diferentes estructuras de recompensa. Asimismo, se
observa que comportamientos como la exploración, explotación, evasión,
coordinación y agrupamiento pueden surgir como consecuencia de la
interacción entre los agentes, el entorno y los incentivos utilizados
durante el aprendizaje.

\end{abstract}

% ---------------------------------------------------------
% 1. INTRODUCCIÓN
% ---------------------------------------------------------

\section{Introducción}

Los ecosistemas pueden entenderse como sistemas dinámicos en los que
múltiples organismos interactúan entre sí y con un ambiente que cambia
constantemente. Las decisiones de cada individuo pueden estar
condicionadas por factores internos, como su nivel energético o estado
fisiológico, y por factores externos, como la disponibilidad de alimento,
la presencia de depredadores, la distribución de recursos y la presencia
de otros individuos. Debido a esta combinación de factores, representar
computacionalmente el comportamiento ecológico constituye un problema que
requiere considerar tanto las características individuales como las
interacciones que ocurren dentro del sistema.

Los Modelos Basados en Agentes (ABM) constituyen una alternativa para
representar estos sistemas mediante individuos autónomos. En un ABM, cada
agente posee un estado y puede ejecutar acciones que modifican su propia
situación o afectan al entorno y a otros agentes. Esta característica ha
permitido utilizar este tipo de modelos para estudiar fenómenos
ecológicos en los que los patrones observados a nivel colectivo son
consecuencia de las decisiones individuales.

Sin embargo, una dificultad de los modelos tradicionales consiste en la
definición de las reglas que controlan el comportamiento de los agentes.
Cuando las decisiones se representan mediante reglas predeterminadas, el
diseñador debe especificar de antemano cómo debe responder el agente ante
las diferentes situaciones posibles. Este enfoque puede resultar
limitante cuando se pretende estudiar escenarios que no fueron
contemplados durante la construcción del modelo.

El Aprendizaje por Refuerzo (RL) proporciona una alternativa basada en el
aprendizaje mediante interacción con el entorno. En lugar de establecer
directamente todas las decisiones mediante reglas, un agente puede
aprender una política que determine qué acción ejecutar a partir del
estado que observa y de las recompensas obtenidas. Esta característica
permite estudiar si determinadas estrategias pueden surgir como
consecuencia del proceso de aprendizaje.

En los trabajos analizados se encuentran aplicaciones de este enfoque a
problemas como el forrajeo, la selección de recursos, la respuesta ante
depredadores, la adaptación a la escasez y la coordinación entre
individuos. Algunos trabajos incorporan además mecanismos de memoria,
variables fisiológicas y observabilidad parcial, buscando representar de
manera más cercana las condiciones bajo las cuales los organismos toman
decisiones.

Un aspecto especialmente relevante para este tipo de modelos es que el
comportamiento aprendido depende de la forma en que se representa el
estado del agente y de la función de recompensa utilizada. Una
representación demasiado limitada puede impedir que el agente disponga
de la información necesaria para tomar una decisión adecuada. Por otro
lado, una recompensa mal definida puede incentivar estrategias que
maximicen la recompensa matemática sin representar el comportamiento
ecológico que se pretendía estudiar.

Este último problema se relaciona con el concepto de
\textit{reward misspecification}, mediante el cual una función de
recompensa no representa de forma adecuada el objetivo que se desea
alcanzar. En sistemas ecológicos, este problema adquiere especial
importancia debido a que conceptos como supervivencia, bienestar,
eficiencia energética o cooperación pueden depender de múltiples
variables y no necesariamente pueden representarse mediante una única
señal de recompensa.

A partir de esta problemática, el presente survey analiza trabajos que
utilizan aprendizaje por refuerzo y modelos basados en agentes para
representar comportamientos individuales y colectivos en ambientes
ecológicos. El análisis se concentra principalmente en la relación entre
las condiciones ambientales, la representación del estado, las
estructuras de recompensa y los comportamientos que emergen durante el
aprendizaje.

La revisión se orienta mediante las siguientes preguntas de
investigación:

\begin{itemize}

    \item \textbf{PI1:} ¿Cómo emergen los comportamientos adaptativos y
    las estrategias de supervivencia colectiva o individual en los
    agentes simulados ante escenarios de estrés ambiental o escasez de
    recursos?

    \item \textbf{PI2:} ¿Cuáles son las variables de estado
    (sensoriales e internas) y las estructuras de funciones de recompensa
    más efectivas para modelar la toma de decisiones ecológicas sin
    generar sesgos o conductas no deseadas
    (\textit{reward misspecification})?

\end{itemize}

El objetivo de la revisión no consiste únicamente en describir los
trabajos seleccionados, sino en identificar patrones comunes entre ellos,
comparar las estrategias utilizadas y determinar qué elementos del diseño
de los agentes parecen ser relevantes para producir comportamientos
adaptativos y ecológicamente plausibles.

% ---------------------------------------------------------
% 2. METODOLOGÍA DE LA REVISIÓN
% ---------------------------------------------------------

\section{Metodología de la revisión}

El presente trabajo corresponde a una revisión de literatura centrada en
el uso de técnicas de aprendizaje por refuerzo y modelos basados en
agentes para representar procesos de decisión y comportamiento ecológico.
El conjunto de análisis está compuesto por 21 trabajos seleccionados por
su relación con las preguntas de investigación planteadas.

Los trabajos abarcan diferentes perspectivas. Algunos estudian
directamente el comportamiento de agentes individuales durante tareas de
forrajeo, mientras que otros analizan sistemas multiagente, coordinación,
competencia, cooperación y formación de estructuras colectivas. También
se incluyen trabajos relacionados con memoria, modelos del mundo,
dinámicas depredador-presa, homeostasis y diseño de entornos de
simulación.

Debido a esta diversidad, la comparación no se realiza únicamente a
partir del algoritmo empleado. En su lugar, los trabajos se analizan
considerando los elementos que determinan el comportamiento del agente y
la manera en que este comportamiento es evaluado.

\subsection{Dimensiones de análisis}

Para realizar la comparación de los trabajos se establecieron cuatro
dimensiones principales:

\begin{enumerate}

    \item \textbf{Representación del estado:} información disponible para
    el agente, incluyendo variables internas, información ambiental,
    información espacial, vecinos y memoria.

    \item \textbf{Función de recompensa:} mecanismo utilizado para
    incentivar determinadas acciones o resultados durante el proceso de
    aprendizaje.

    \item \textbf{Entorno e interacción:} características del ambiente de
    simulación, disponibilidad de recursos, amenazas, número de agentes
    y nivel de observabilidad.

    \item \textbf{Comportamiento y evaluación:} estrategias individuales
    o colectivas observadas y métricas utilizadas para determinar si el
    comportamiento aprendido cumple con el objetivo planteado.

\end{enumerate}

Estas dimensiones permiten comparar estudios que emplean algoritmos
diferentes pero que presentan problemas ecológicos similares. Por
ejemplo, un trabajo puede utilizar aprendizaje por refuerzo profundo para
estudiar forrajeo individual, mientras otro puede utilizar aprendizaje
multiagente para estudiar coordinación. Aunque los algoritmos sean
distintos, ambos pueden proporcionar evidencia sobre la influencia de la
representación del estado o de la función de recompensa en el
comportamiento obtenido.

\subsection{Clasificación temática}

A partir de las dimensiones anteriores, los trabajos se organizarán en
cuatro áreas temáticas principales:

\begin{enumerate}

    \item \textbf{Forrajeo y adaptación ante escasez de recursos.}
    Incluye trabajos relacionados con búsqueda de alimento, explotación
    de recursos, exploración, agotamiento de parches y supervivencia.

    \item \textbf{Estado, percepción y memoria.}
    Comprende investigaciones que analizan variables fisiológicas,
    información espacial, observabilidad parcial, memoria y
    representaciones internas.

    \item \textbf{Recompensas y homeostasis.}
    Incluye trabajos centrados en funciones de recompensa, objetivos
    individuales o colectivos, homeostasis y problemas derivados de una
    especificación incorrecta de las recompensas.

    \item \textbf{Comportamiento colectivo y sistemas multiagente.}
    Comprende estudios sobre cooperación, competencia, coordinación,
    agrupamiento, formación de bandadas y otros comportamientos
    colectivos.

\end{enumerate}

Esta clasificación se utilizará posteriormente para organizar la
discusión de los resultados. Un mismo trabajo puede aportar evidencia a
más de una categoría, debido a que la representación del estado, la
recompensa y la interacción entre agentes se encuentran estrechamente
relacionadas.

\subsection{Criterios de comparación}

Además de la clasificación temática, cada trabajo será comparado
considerando las siguientes características:

\begin{table}[H]
\centering
\caption{Dimensiones utilizadas para el análisis comparativo}
\label{tab:dimensiones}
\small
\begin{tabular}{p{3.5cm} p{4.5cm}}
\toprule
\textbf{Dimensión} & \textbf{Aspecto considerado} \\
\midrule

Algoritmo &
Técnica de aprendizaje por refuerzo utilizada y paradigma de
entrenamiento. \\

Estado &
Variables internas, ambientales, espaciales, sociales y temporales
disponibles para el agente. \\

Observabilidad &
Información completa o parcial disponible durante la toma de
decisiones. \\

Recompensa &
Tipo de recompensa, objetivo incentivado y nivel de especificidad. \\

Entorno &
Características del ambiente, distribución de recursos, amenazas y
dinámica ambiental. \\

Interacción &
Número de agentes y presencia de cooperación, competencia o
coordinación. \\

Comportamiento &
Estrategias individuales o colectivas observadas durante la simulación. \\

Evaluación &
Métricas utilizadas para evaluar aprendizaje, supervivencia y
comportamiento emergente. \\

\bottomrule
\end{tabular}
\end{table}

% ---------------------------------------------------------
% 3. FUNDAMENTOS CONCEPTUALES
% ---------------------------------------------------------

\section{Fundamentos conceptuales}

Esta sección presenta los conceptos necesarios para interpretar los
resultados de la revisión. Debido a que el objetivo principal del trabajo
es analizar la literatura y no desarrollar una introducción extensa a
cada técnica, los conceptos se presentan únicamente desde la perspectiva
de su utilización en sistemas ecológicos.

\subsection{Modelos Basados en Agentes}

Los Modelos Basados en Agentes permiten representar sistemas complejos
mediante individuos autónomos que interactúan entre sí y con un entorno.
Cada agente puede poseer un conjunto de variables que describe su estado,
un conjunto de acciones disponibles y mecanismos para modificar su
comportamiento.

En un contexto ecológico, un agente puede representar un organismo que
posee características como posición, energía, nivel de hambre,
información sensorial y experiencias anteriores. El entorno puede incluir
recursos, obstáculos, depredadores y otros individuos.

Una característica importante de los ABM es que permiten estudiar cómo
las decisiones individuales pueden producir patrones a nivel colectivo.
Por esta razón, resultan adecuados para investigar fenómenos como
agrupamiento, cooperación, competencia o formación de bandadas.

Sin embargo, en un ABM tradicional las decisiones pueden depender de
reglas diseñadas manualmente. El aprendizaje por refuerzo permite
modificar este enfoque al permitir que el agente aprenda una política de
decisión mediante la interacción con el entorno.

\subsection{Aprendizaje por Refuerzo}

El Aprendizaje por Refuerzo estudia el proceso mediante el cual un agente
aprende a seleccionar acciones a partir de su interacción con un entorno.
En una formulación simplificada, el proceso puede representarse como:

\begin{equation}
s_t \rightarrow a_t \rightarrow r_t \rightarrow s_{t+1}
\end{equation}

donde $s_t$ representa el estado del agente en el instante $t$, $a_t$
representa la acción seleccionada, $r_t$ corresponde a la recompensa
obtenida y $s_{t+1}$ representa el nuevo estado.

El objetivo del agente consiste en aprender una política que permita
obtener recompensas acumuladas elevadas. En aplicaciones ecológicas, las
acciones pueden representar desplazarse, consumir alimento, descansar,
huir de un depredador o interactuar con otros individuos.

La importancia de este enfoque para la presente revisión radica en que
permite estudiar qué estrategias aparecen cuando el agente debe aprender
a partir de las consecuencias de sus propias acciones.

\subsection{Aprendizaje por Refuerzo Profundo}

El Aprendizaje por Refuerzo Profundo (DRL) combina el aprendizaje por
refuerzo con redes neuronales utilizadas para aproximar funciones de
valor, políticas o representaciones del estado.

Su utilización resulta relevante cuando el espacio de estados presenta
una cantidad elevada de variables o cuando las relaciones entre las
variables no pueden representarse fácilmente mediante estructuras
tabulares.

Dentro de los trabajos revisados, el DRL aparece asociado a problemas
como forrajeo, conservación, percepción visual y modelado de ecosistemas.
Sin embargo, el uso de una red neuronal por sí mismo no garantiza la
aparición de un comportamiento ecológicamente válido. La información
proporcionada al agente y la función de recompensa continúan siendo
elementos fundamentales.

\subsection{Observabilidad parcial y memoria}

En un ambiente ecológico, un agente normalmente no tiene acceso a toda la
información existente en el entorno. Su percepción puede estar limitada
por la distancia, el campo de visión, obstáculos o las capacidades
sensoriales del organismo representado.

Esta situación puede formularse mediante un Proceso de Decisión de Markov
Parcialmente Observable (POMDP), donde el agente debe tomar decisiones
utilizando únicamente una observación parcial del estado real del
entorno.

La observabilidad parcial hace que la memoria pueda adquirir un papel
importante. Un agente puede utilizar información de experiencias
anteriores para complementar su percepción actual. Esta información puede
corresponder a la ubicación de recursos, el tiempo desde una visita
anterior, la calidad de una zona o experiencias relacionadas con una
amenaza.

Por lo tanto, la memoria puede considerarse un componente de la
representación del estado cuando la información necesaria para tomar una
decisión no está disponible directamente en el momento actual.

\subsection{Aprendizaje por Refuerzo Multi-Agente}

Cuando varios agentes aprenden e interactúan simultáneamente, el problema
se extiende al Aprendizaje por Refuerzo Multi-Agente (MARL). En este caso,
las acciones de un agente pueden modificar el entorno observado por los
demás.

En un escenario ecológico, por ejemplo, varios individuos pueden consumir
el mismo recurso, competir por una fuente de alimento o beneficiarse de
permanecer juntos. Como consecuencia, el aprendizaje de cada agente puede
afectar el proceso de aprendizaje de los demás.

Esta interacción introduce problemas como la no-estacionariedad,
coordinación y asignación de crédito (\textit{credit assignment}). Al
mismo tiempo, permite estudiar cómo pueden surgir comportamientos
colectivos a partir de decisiones individuales.

\subsection{Comportamiento emergente}

Para efectos de esta revisión, se considera comportamiento emergente
aquel patrón observable que resulta de la interacción entre los agentes,
el entorno y las reglas de aprendizaje, sin necesidad de definir
explícitamente dicho patrón como una regla de comportamiento.

Entre los comportamientos analizados en los trabajos se encuentran la
cooperación, competencia, coordinación, formación de grupos, evasión
colectiva y adaptación ante cambios en la disponibilidad de recursos.

La distinción entre comportamiento programado y comportamiento emergente
es importante para esta revisión. Si una función de recompensa obliga
directamente al agente a producir un patrón determinado, observar dicho
patrón no necesariamente demuestra que haya surgido como consecuencia de
la interacción con el entorno.

Por esta razón, los trabajos se analizarán no solamente considerando qué
comportamiento aparece, sino también las condiciones bajo las cuales
aparece y la forma en que los autores determinan que dicho comportamiento
es realmente emergente.

% ---------------------------------------------------------
% 4. Análisis temático de la literatura
% ---------------------------------------------------------

\section{Análisis de la literatura}

Como paso previo al desglose conceptual, y para facilitar la contrastación 
metodológica de los distintos enfoques, la 
\textbf{Tabla \ref{tab:comparativa}} (ubicada en esta sección) sintetiza los 
trabajos más representativos de la revisión. Esta matriz cruza las 
técnicas algorítmicas, los entornos de simulación (datasets/Gridworlds), 
las métricas principales y, de manera crítica, las limitaciones 
identificadas en cada estudio. A partir del mapeo expuesto en la tabla, 
la literatura se analiza a través de las siguientes cuatro dimensiones 
fundamentales:

\subsection{Forrajeo y adaptación ante la escasez de recursos}
La literatura demuestra que el forrajeo óptimo ante la escasez no depende de 
reglas fijas, sino de la capacidad de anticipación del agente. Los agentes que 
construyen representaciones internas o ''modelos del mundo'' logran planificar 
trayectorias de forrajeo óptimas ante recursos heterogéneos, convergiendo a 
estrategias de la teoría clásica del Teorema del Valor Marginal sin requerir 
memorias tabulares rígidas. Asimismo, el agotamiento gradual de los recursos 
actúa como una fuente crucial de información; las aves aprenden la calidad 
inicial del parche basándose en el patrón de espaciado temporal entre bocados, 
optimizando su supervivencia mediante reglas de temporización en entornos 
variables. Además, los algoritmos de aprendizaje sin modelo (model-free) han 
demostrado descubrir estrategias autónomas ante situaciones críticas y puntos 
de inflexión ecológica, superando las reglas empíricas humanas frente a 
perturbaciones ambientales.

\subsection{Representación del estado, percepción y memoria}
Una representación de estado efectiva requiere combinar la percepción espacial 
con variables fisiológicas. La inclusión de señales de interocepción, como 
el nivel de saciedad, permite a los agentes adaptar su velocidad y realizar 
pausas voluntarias para conservar energía \cite{Paper1}. En entornos de 
escasez extrema, incorporar mecanismos de memoria episódica biológica resulta 
vital para que el agente recuerde qué, dónde y cuándo almacenó comida, 
evitando visitar celdas recién agotadas. El estado también suele estructurarse 
combinando variables internas de salud con información de una vecindad espacial
limitada ($3 x 3$) y señales de distancia a los recursos \cite{Paper2}. 
Para garantizar el realismo de estas simulaciones, los modelos deben basarse 
en resoluciones espaciales precisas que capturen la dinámica de la vegetación 
y el agua, elementos fundamentales para determinar la idoneidad del hábitat.

\subsection{IV-C. Funciones de recompensa y homeostasis}
El diseño de las recompensas enfrenta el riesgo de generar conductas sesgadas 
(reward misspecification), donde el agente aprende atajos matemáticos en lugar 
de estrategias ecológicas. Para evitar recompensas dispersas, se proponen 
funciones de recompensa homeostáticas cuadráticas que incentivan mantener 
niveles internos de energía e hidratación equilibrados \cite{Paper2}. Estas 
funciones no lineales penalizan cuadráticamente tanto la deficiencia como la 
sobresaturación, obligando a la IA a buscar un equilibrio biológico en lugar 
de maximizar el consumo infinitamente. Adicionalmente, el estado interno debe 
interactuar con las amenazas del entorno; la ponderación de la recompensa 
mediante penalizaciones por riesgo de depredación permite modelar el conflicto 
biológico entre el hambre y el miedo. En simulaciones multiagente con escasez, 
las recompensas diseñadas a mano suelen fallar, haciendo necesaria la 
integración de redes predictoras (MARP) que evalúen el éxito a nivel 
poblacional en lugar del beneficio individual.

\subsection{IV-D. Comportamiento colectivo y sistemas multiagente}
En entornos multiagente, el comportamiento emergente surge sin necesidad de 
programar reglas rígidas. En cuadrículas 2D acotadas, agentes perseguidores 
pueden desarrollar hasta 21 patrones de coordinación táctica excluyendo 
deliberadamente las recompensas densas por distancia y utilizando únicamente 
recompensas de resultado. La formación de bandadas con alto orden polar emerge 
de manera natural al aplicar funciones de costo local que penalizan tanto la 
dispersión excesiva como el hacinamiento entre individuos. Incluso en 
escenarios donde los individuos son egoístas y carecen de incentivos 
explícitos para cooperar, la agrupación (swarming) surge como estrategia 
adaptativa para compensar la observabilidad parcial del entorno. Sin embargo, 
la literatura advierte que un alto desempeño en términos de recompensa media 
no siempre implica cooperación biológica real, por lo que es necesario 
analizar estadísticamente las trayectorias para validar si la colaboración es 
causal o una ilusión estadística.

\subsection{Comparación de estrategias de modelado}

La comparación de los trabajos muestra que la representación del estado
constituye uno de los principales factores que diferencian los enfoques
revisados. Los modelos más simples utilizan principalmente información
espacial o distancia hacia recursos y otros agentes, mientras que los
modelos orientados a representar procesos ecológicos más complejos
incorporan variables internas como energía, hambre, hidratación o
saciedad.

También se observa una diferencia entre los modelos que utilizan
información directamente observable y aquellos que incorporan memoria o
representaciones internas. En ambientes donde los recursos cambian o
presentan una distribución temporal, la memoria permite utilizar
información obtenida en interacciones anteriores. De manera similar, la
observabilidad parcial obliga al agente a tomar decisiones sin disponer
del estado completo del ambiente.

En cuanto a las funciones de recompensa, los trabajos revisados
presentan desde recompensas simples asociadas con la obtención de
recursos hasta estructuras que consideran homeostasis, supervivencia,
riesgo y objetivos colectivos. Esta diferencia resulta especialmente
importante debido a que una recompensa demasiado específica puede
restringir el comportamiento del agente, mientras que una recompensa
demasiado general puede permitir estrategias que maximicen la señal
numérica sin representar adecuadamente el objetivo ecológico.

En los sistemas multiagente, esta situación se vuelve más compleja debido
a que el comportamiento de un agente puede modificar las condiciones
experimentadas por los demás. Por ello, algunos estudios utilizan
recompensas individuales, mientras que otros incorporan objetivos
compartidos o multiobjetivo. Los resultados revisados sugieren que la
cooperación y la coordinación no necesariamente deben programarse de
forma explícita, pero su aparición debe evaluarse cuidadosamente para
distinguir entre un comportamiento verdaderamente emergente y uno
producido directamente por el diseño de la recompensa.

\section{Síntesis de los hallazgos}

La literatura revisada permite observar que el comportamiento adaptativo
de los agentes no depende exclusivamente del algoritmo de aprendizaje
utilizado. La interacción entre la representación del estado, la función
de recompensa y las características del entorno tiene un papel
fundamental en las estrategias obtenidas.

En escenarios de escasez de recursos, los agentes pueden modificar sus
estrategias de exploración y explotación, utilizar información obtenida
de experiencias anteriores y responder a cambios en la disponibilidad de
recursos. Cuando además se incorporan múltiples agentes, estas decisiones
individuales pueden generar patrones colectivos como agrupamiento,
coordinación, cooperación o competencia.

De manera general, los estudios también muestran que incrementar la
complejidad del modelo no garantiza un comportamiento más realista. La
incorporación de variables adicionales debe responder a las capacidades
que se desean representar y a la información que tendría disponible el
organismo modelado. De forma similar, una función de recompensa más
compleja no necesariamente produce mejores resultados si introduce
incentivos contradictorios o comportamientos no deseados.

Estos hallazgos establecen la base para analizar posteriormente las
preguntas de investigación. En particular, permiten estudiar qué
combinaciones de variables de estado y recompensas favorecen la
adaptación individual y cuáles permiten que aparezcan comportamientos
colectivos sin definirlos explícitamente.

% --- TABLA COMPARATIVA DE LOS ESTUDIOS ---
\begin{table*}[t]
\centering
\caption{Comparación de los trabajos revisados}
\label{tab:comparativa}
\scriptsize
\begin{tabular}{p{2.4cm} p{1.2cm} p{1.7cm} p{2.5cm} p{2.7cm} p{3.0cm}}
\toprule
\textbf{Trabajo} &
\textbf{Año} &
\textbf{Enfoque} &
\textbf{Estado / Variables} &
\textbf{Recompensa / Objetivo} &
\textbf{Principal aporte} \\
\midrule

Jiang et al. \cite{Paper1} &
2024 &
DRL / ecología visual &
Alimento, señales visuales, saciedad y estado interno &
Supervivencia y obtención de alimento &
Los agentes aprenden a distinguir recursos beneficiosos y perjudiciales. La memoria recurrente y el estado de saciedad permiten desarrollar estrategias de forrajeo más complejas. \\

\addlinespace

Agent-based ecosystem modeling \cite{Paper2} &
2024 &
DRL + ABM &
Energía, hidratación, vecinos, recursos y obstáculos &
Recompensa homeostática &
Muestra que ciclos depredador-presa y patrones espaciales pueden emerger a partir de agentes con reglas de decisión aprendidas. \\

\addlinespace

Deep reinforcement learning for conservation decisions \cite{Paper3} &
2024 &
DRL / TD3 &
Estado ambiental y condiciones de perturbación &
Supervivencia y conservación &
Los agentes pueden descubrir estrategias ante puntos de inflexión ecológica, aunque la formulación de la recompensa y del POMDP constituye una fuente de riesgo. \\

\addlinespace

World Models \cite{Paper4} &
2025 &
RL + World Models &
Representación latente del entorno y recursos &
Energía neta / eficiencia de forrajeo &
Las representaciones predictivas del entorno permiten anticipar cambios y desarrollar estrategias de forrajeo cercanas a predicciones de la teoría del forrajeo óptimo. \\

\addlinespace

Episodic-like memory \cite{Paper5} &
2023 &
Modelo cognitivo &
Qué, dónde y cuándo ocurrió un evento &
Supervivencia / recuperación de alimento &
Resalta la importancia de la memoria espacial y temporal para tomar decisiones durante períodos de escasez. \\

\addlinespace

Cognitive-Emotional Forager \cite{Paper6} &
2005 &
Red neuronal / forrajeo &
Hambre, energía, amenaza y miedo &
Equilibrio entre alimentación y riesgo &
Modela el conflicto entre satisfacer necesidades internas y evitar depredadores. \\

\addlinespace

Group foraging behaviors \cite{Paper7} &
--- &
Modelo de comportamiento grupal &
Posición y densidad de vecinos, movimiento &
Éxito individual / competencia &
Relaciona decisiones individuales con patrones colectivos de movimiento durante el forrajeo. \\

\addlinespace

Decision-Making in ABM \cite{Survey1} &
--- &
Survey &
Estado, recompensa, interacción y decisión &
Diversos &
Identifica estructuras recurrentes para representar decisiones en modelos basados en agentes y destaca la importancia de los modelos del mundo y la interacción. \\

\addlinespace

Evolutionary game dynamics \cite{Paper9} &
2026 &
RL + teoría de juegos &
Estado estratégico e interacción &
Cooperación / competencia &
Relaciona el aprendizaje por refuerzo con procesos de adaptación, competencia y coordinación entre agentes. \\

\addlinespace

AI-enhanced ABM for wild populations \cite{Paper10} &
2025 &
IA + ABM &
Condiciones ambientales y poblacionales &
Supervivencia poblacional &
Propone modelos capaces de representar poblaciones adaptativas ante perturbaciones como sequías y cambios del paisaje. \\

\addlinespace

Reward Prediction / MARP \cite{Paper11} &
--- &
MARL &
Estado individual y colectivo &
Objetivos individuales y globales &
Plantea mecanismos para alinear las recompensas individuales con objetivos colectivos y reducir comportamientos no deseados. \\

\addlinespace

Emergent Coordination \cite{Paper12} &
2025 &
MARL independiente &
Identidad, densidad y configuración espacial &
Coordinación colectiva &
Muestra transiciones entre comportamientos coordinados y no coordinados bajo diferentes condiciones de densidad y congestión. \\

\addlinespace

Resource depletion \cite{Paper13} &
2026 &
RL / forrajeo &
Cantidad, composición y disponibilidad temporal &
Obtención eficiente de recursos &
El agotamiento de los recursos funciona como información que permite modificar la estrategia de exploración y explotación. \\

\addlinespace

Homeostasis and resource sharing \cite{Paper14} &
2025 &
Multi-objective MARL &
Energía, recursos y estado interno &
Homeostasis y sostenibilidad &
Integra objetivos biológicos y colectivos para estudiar el consumo y distribución sostenible de recursos. \\

\addlinespace

Collaborative emergent behavior \cite{Paper15} &
2018 &
MADDPG / MARL &
Trayectorias e interacción entre agentes &
Captura / recompensa colectiva &
Advierte que una recompensa elevada no demuestra necesariamente cooperación real, por lo que se requiere analizar las trayectorias y relaciones causales. \\

\addlinespace

Pursuit-Evasion Games \cite{Paper16} &
2025 &
MARL &
Posición y movimiento en Gridworld &
Resultado de persecución &
Identifica múltiples patrones tácticos y comportamientos coordinados emergentes en escenarios de persecución y evasión. \\

\addlinespace

Learning to Flock \cite{Paper17} &
2026 &
Q-learning &
Distancia relativa entre vecinos &
Cohesión y evitación de colisiones &
La formación de bandadas puede emerger utilizando reglas de aprendizaje simples basadas en mantener proximidad sin producir colisiones. \\

\addlinespace

Foraging Swarms \cite{Paper18} &
2020 &
MARL &
Observaciones locales y posición de agentes &
Obtención de recursos &
Muestra que el agrupamiento puede surgir incluso cuando los agentes reciben recompensas individuales, debido a los beneficios de la observación colectiva. \\

\addlinespace

Aquarium \cite{Paper19} &
2024 &
MARL / PPO &
Campo de visión, vecinos, recursos y amenazas &
Supervivencia / reproducción &
Proporciona un entorno multiagente para estudiar dinámicas depredador-presa, reproducción, inanición y comportamientos colectivos. \\

\addlinespace

Development of Swarm Behavior \cite{Paper20} &
2020 &
Projective Simulation &
Posición y distancia a recursos &
Forrajeo y agrupamiento &
Relaciona la distribución de recursos con la aparición de comportamiento colectivo y diferentes patrones de movimiento. \\

\addlinespace

Multi-Agent Deep Reinforcement Learning Survey \cite{Survey3} &
2021 &
Survey de MARL &
Estados individuales y globales &
Cooperación / competencia &
Resume problemas fundamentales de MARL, incluyendo no estacionariedad, asignación de crédito, escalabilidad y diseño de recompensas. \\

\bottomrule
\end{tabular}
\end{table*}

La comparación presentada en la Tabla \ref{tab:comparativa} permite
identificar que los trabajos revisados no utilizan un único enfoque para
representar el comportamiento ecológico. Por el contrario, existe una
relación recurrente entre tres elementos: la información disponible para
el agente, la función utilizada para evaluar sus decisiones y las
interacciones que ocurren dentro del entorno.

A partir de esta relación, los estudios pueden agruparse principalmente
en cuatro dimensiones: adaptación durante el forrajeo y la escasez de
recursos, representación del estado y mecanismos de percepción y memoria,
diseño de funciones de recompensa y homeostasis, y comportamiento
colectivo en sistemas multiagente. Estas dimensiones permiten analizar los
trabajos de forma transversal y relacionarlos directamente con las
preguntas de investigación planteadas.

% ---------------------------------------------------------
% REFERENCIAS
% ---------------------------------------------------------

\begin{thebibliography}{99}

\bibitem{Paper1}
A computational approach to visual ecology with deep reinforcement
learning.

\bibitem{Paper2}
Agent-based ecosystem modeling with deep reinforcement learning.

\bibitem{Paper3}
Deep reinforcement learning for conservation decisions.

\bibitem{Paper4}
World Models Unlock Optimal Foraging Strategies in Reinforcement
Learning Agents.

\bibitem{Paper5}
Computational models of episodic-like memory in food-caching birds.

\bibitem{Paper6}
A neural network model of foraging decisions made under predation risk:
The Cognitive-Emotional Forager.

\bibitem{Paper7}
Linking cognitive strategy, neural mechanism, and movement statistics in
group foraging behaviors.

\bibitem{Paper8}
A review on vision-centric coarse to fine-grained animal action
recognition.

\bibitem{Survey1}
Decision-Making in Agent-Based Modeling: A Current Review and Future
Prospectus.

\bibitem{Paper9}
A brief review of evolutionary game dynamics in the reinforcement
learning paradigm.

\bibitem{Paper10}
Using AI enhanced agent-based models to support management of wild
populations.

\bibitem{Paper11}
Why Study Emergent Behavior When You Can Regulate It? Aligning
Multi-Agent Systems with Reward Prediction.

\bibitem{Paper12}
Emergent Coordination and Phase Structure in Independent Multi-Agent
Reinforcement Learning.

\bibitem{Paper13}
Resource depletion accelerates rate learning but not composition learning
in patch foraging.

\bibitem{Paper14}
From homeostasis to resource sharing: Biologically and economically
aligned multi-objective multi-agent gridworld-based AI safety
benchmarks.

\bibitem{Survey2}
Machine and Deep Learning for Wetland Mapping and Bird-Habitat Monitoring: A Systematic Review of Remote-Sensing Applications

\bibitem{Paper15}
Measuring collaborative emergent behavior in multiagent reinforcement
learning.

\bibitem{Paper16}
Emergent Behaviors in Multiagent Pursuit-Evasion Games Within a Bounded
2D Grid World.

\bibitem{Paper17}
Learning to Flock in Open Space by Avoiding Collisions and Staying
Together.

\bibitem{Paper18}
Foraging Swarms Using Multi-Agent Reinforcement Learning.

\bibitem{Paper19}
Aquarium: A Comprehensive Framework for Exploring Predator-Prey
Dynamics through Multi-Agent Reinforcement Learning Algorithms.

\bibitem{Paper20}
Development of Swarm Behavior in Artificial Learning Agents That Adapt
to Different Foraging Environment.

\bibitem{Paper22}
World Models Unlock Optimal Foraging Strategies in Reinforcement Learning Agents

\bibitem{Survey3}
Multi-Agent Deep Reinforcement Learning: A Survey.

\end{thebibliography}



\end{document}

